\section{Relative affine-branch closure}
\label{sec:theorem}

For $r\ge2$, write
\[
M_r=\langle u,v\mid vu=u^rv\rangle^+,
\qquad
G_r=\langle u,v\mid vuv^{-1}=u^r\rangle.
\]
The predecessor objects associated with these presentations are not identified
with each other.  The symbols only index the source-owned families whose
outcomes are recorded at the expanded endpoints.

\subsection{Endpoint-obstruction totality}

Let $V_{\mathrm{test}}$ be the 17 tested endpoints in the expanded graph.  The
bridge stores the following total map to failed coordinates:
\begin{align*}
\Fail(\code{N35F})=\Fail(\code{N35P})&=\{R,D\},\\
\Fail(\code{N35S})=\Fail(\code{N35H})&=\{S,D,C\},\\
\Fail(\code{N35Q})&=\{S,D,M\},\\
\Fail(\code{N35D})&=\{R,S,M\},\\
\Fail(\code{N35B})&=\{I,S,M\},\\
\Fail(\code{N36F})&=\{R,S,D,M,C\},\\
\Fail(\code{N36G})&=\{R,S,C\},\\
\Fail(\code{N37O})=\Fail(\code{N37D})&=\{S,C\},\\
\Fail(\code{N37N})&=\{R,S,C\},\\
\Fail(\code{N38T})=\Fail(\code{N38M})&=\{R,D\},\\
\Fail(\code{N38L})&=\{D,C\},\\
\Fail(\code{N38O})&=\{S,D,M,C\},\\
\Fail(\code{N38K})&=\{M\}.
\end{align*}
Every value is nonempty.  The map is not defined by graph reachability alone;
each set is inherited from the source-owned theorem at that endpoint.  The
sink \code{NX} instead has an empty failed-coordinate set because it represents
nonmembership, not an unsuccessful candidate.

\begin{lemma}[Endpoint-obstruction totality]
\label{lem:endpoint-totality}
For every $p\in\Paths$, the endpoint $\End(p)$ lies in
$V_{\mathrm{test}}$.  For every $c\in X_{\End(p)}(r)$, at least one coordinate
in $\Fail(\End(p))$ is false for $c$.
\end{lemma}

\begin{proof}
The finite path table lists one nonempty provenance path for each tested
endpoint and no in-contract path to \code{NX}, \code{NC}, \code{NR}, or
\code{NS}.  Direct inspection of the edge endpoints shows that the terminal
node of every listed in-contract path is one of the 17 nodes above.  The
inherited theorem named by the endpoint record proves the stated coordinate
failure under that endpoint's own object and quantifiers.  The nonempty sets
displayed above exhaust $V_{\mathrm{test}}$, which proves both assertions.
\end{proof}

The word ``totality'' in \cref{lem:endpoint-totality} is deliberately local.
It refers to the finite endpoint table.  It is not a claim that all conceivable
repair processes terminate, nor does it use reflexive reachability to manufacture
a terminal witness.

\subsection{The token classifier}

The normalizer in \cref{eq:normalizer} assigns eight tokens to obstruction
paths and eight to exit paths.  The obstruction tokens are the six frozen
families plus the canonical Bass--Serre and groupoid-import tokens.  The exit
tokens are the six exit-only requests plus the two explicit alternatives from
the mixed classes.  Every token record names its exact instance scope; the
normalizer cannot be extended by imagining an unrecorded instance.

\begin{proposition}[Finite request classification]
\label{prop:token-totality}
For every $q\in\SigmaReq$, exactly one of the following holds.
\begin{enumerate}
  \item $\mathsf{norm}(q)$ is an obstruction record with at least one endpoint
  $v\in V_{\mathrm{test}}$ and $\Fail(v)\ne\varnothing$.
  \item $\mathsf{norm}(q)$ is an exit record whose terminal endpoint is
  \code{NX}, whose obstruction-endpoint list is empty, and whose failed-\Good{}
  set is empty.
\end{enumerate}
\end{proposition}

\begin{proof}
The 16 token IDs in \cref{tab:tokens} are distinct and equal the declared
domain of the finite normalizer.  Eight records have disposition
\Obstructed{} and nonempty obstruction-endpoint/path lists; their endpoints
belong to $V_{\mathrm{test}}$, so \cref{lem:endpoint-totality} applies.  The
remaining eight records have disposition \Exit{}, empty obstruction lists, and
an explicit exit edge ending at \code{NX}.  No record has both dispositions,
and the two counts sum to 16.
\end{proof}

\Cref{fig:classification} visualizes this finite, disjoint classifier.

\begin{figure}[H]
  \centering
  \resizebox{0.96\textwidth}{!}{\begin{tikzpicture}[
  x=1cm,y=1cm,>=Latex,
  classbox/.style={rounded corners=2pt,draw=charcoal,very thick,fill=white,
    align=left,text width=4.35cm,minimum height=4.15cm,font=\scriptsize,
    inner sep=5pt},
  obstructed/.style={classbox,draw=deepgreen},
  mixed/.style={classbox,draw=formalblue,double},
  exitonly/.style={classbox,draw=warningamber,dashed},
  endpoint/.style={rounded corners=2pt,draw=charcoal,thick,align=center,
    text width=4.2cm,minimum height=.95cm,font=\small,fill=white},
  solidarr/.style={->,very thick,deepgreen},
  exitarr/.style={->,very thick,dashed,warningamber},
  note/.style={rounded corners,draw=charcoal,fill=softgray,align=center,
    font=\scriptsize,inner sep=4pt}
]
\node[obstructed] (obs) at (0,1.3) {\textbf{6 obstructed classes}\\[2pt]
  affine Cayley representation\\
  finite-rank local system\\
  character\\grading\\quotient\\modular phase\\[4pt]
  \textbf{6 frozen-family tokens}\\
  each has an in-domain witness path};

\node[mixed] (mix) at (5.45,1.3) {\textbf{2 mixed classes}\\[2pt]
  Bass--Serre splitting\\
  groupoid trace\\[6pt]
  \textbf{2 canonical tokens}\\
  frozen splitting / frozen groupoid import\\[4pt]
  \textbf{2 alternative EXIT tokens}\\
  changed splitting / determinant category};

\node[exitonly] (ext) at (10.9,1.3) {\textbf{6 exit-only classes}\\[2pt]
  induced shift\\first-return map\\valuation tree\\
  boundary model\\basepoint damping\\finite-total-weight retrofit\\[4pt]
  \textbf{6 explicit EXIT tokens}\\
  no canonical obstruction credit};

\node[endpoint,draw=deepgreen,below=9mm of obs] (bad)
  {\textbf{obstruction endpoint}\\nonempty failed-\texttt{Good} set};
\node[endpoint,draw=warningamber,dashed,below=9mm of ext] (out)
  {\textbf{typed EXIT endpoint}\\empty failed-\texttt{Good} set};

\draw[solidarr] (obs.south) -- (bad.north);
\draw[solidarr] (mix.south west) -- (bad.north east);
\draw[exitarr] (mix.south east) -- (out.north west);
\draw[exitarr] (ext.south) -- (out.north);

\node[note,below=6mm of mix,text width=4.55cm] (total)
  {$8$ obstructed tokens $+8$ EXIT tokens $=\Sigma_{16}$\\
   no \texttt{OTHER\_INSTANCE}; no implicit compound request};
\draw[solidarr] (bad.east) -- (total.west);
\draw[exitarr] (out.west) -- (total.east);

\node[note,below=4mm of total,text width=10.2cm,fill=white] (scope)
  {class dispositions are not token outcomes: mixed classes contribute one
   obstructed canonical token and one explicit category-changing EXIT token};
\end{tikzpicture}
}
  \caption{The finite classifier separates class dispositions from token
  outcomes.  Six classes are obstructed and six are exit-only.  Each of the two
  mixed classes contributes one obstructed canonical token and one explicit
  category-changing exit token.  Thus the token census is eight plus eight.
  Solid paths terminate with a nonempty failed-\Good{} set; dashed EXIT paths
  terminate with an empty set and never count as evidence for $\neg\Good$.}
  \label{fig:classification}
\end{figure}

At class level, the census is six obstructed, six exit-only, and two mixed
(Bass--Serre splitting and groupoid trace).  An exit-only row says only that its
one request token changes the frozen category; a mixed row quantifies only over
its two recorded tokens, not all conceivable alternatives.

\subsection{Relative closure theorem}

\begin{theorem}[Relative exhaustion of the encoded affine branch]
\label{thm:closure}
Under \cref{ass:timing,ass:predecessors}, for every $r\ge2$ and every
$c\in\Caff(r)$,
\[
\neg\Good(c).
\]
Moreover, every request $q\in\SigmaReq$ is classified by
\cref{prop:token-totality}; no request in the finite domain remains unassigned.
\end{theorem}

\begin{proof}
Take $c\in\Caff(r)$.  By \cref{eq:candidate-domain}, $c=(p,x)$ for a nonempty
provenance path $p\in\Paths$ and an endpoint datum
$x\in X_{\End(p)}(r)$.  By \cref{lem:endpoint-totality}, the endpoint lies in
$V_{\mathrm{test}}$ and at least one coordinate listed by
$\Fail(\End(p))$ is false for $x$.  Since \Good{} is the conjunction of all
six coordinates, $\Good(c)$ is false.  The final request-classification
statement is \cref{prop:token-totality}.
\end{proof}

The theorem has two related conclusions but no hidden equivalence.  The first
quantifies over endpoint-typed candidate data on the in-contract provenance
paths.  The second classifies all 16 request tokens, including category exits
that do not lie in $\Caff(r)$.  Exit classification completes the request table;
it does not prove $\neg\Good$ for a changed object.

\subsection{Inherited kernels and sharp boundaries}

The endpoint map packages four source-owned kernels: Paper 35 supplies the
object fork; Paper 36 the filling, quotient-ownership, and clock fork; Paper 37
the finite-coefficient leakage-or-erasure fork; and Paper 38 the tree/orbital
trilemma for the tree canonical to the frozen ascending-HNN splitting and
presentation.  Paper 39 contributes their typed coverage proof, not new
versions of these kernels.  Their hypotheses and dependency structure remain
in the frozen packages and \cref{app:ledgers}.

The quantifier $r\ge2$ is intentional.  At $r=1$ some affine obstructions
disappear while filling, nonproper-action, and orbital controls remain, so the
balanced case is only a control.  The explicit countermodels in
\cref{app:ledgers}---finite cycles, noninvertible deletion, summable tree
damping, proper tree lattices, and first-trace cancellation---show sharply why
\cref{thm:closure} cannot become a universal impossibility theorem.
